\subsection{Deducing forbidden outside paths}
\label{finite:constraints}

For each pattern, we derive restrictions that every graph realizing the
chosen optimal system must satisfy.  These restrictions will supply
separated balls for the independence argument.

Let \(X\) be a vertex set containing \(S\).  For distinct \(x,y\in X\),
an \emph{outside path relative to \(X\) with parameter \(d\)} is a simple
\(x\)-\(y\) path with exactly \(d\) internal vertices, all in
\(V(G)\setminus X\).  Here \(d\) is a nonnegative integer; in particular,
parameter \(0\) means the edge \(xy\).  We maintain sets
\[
 M_{xy}\subseteq\{0,1,2,\ldots\},\qquad
 M_{yx}=M_{xy},\qquad M_{xx}=\varnothing,
\]
with the interpretation that no such path exists when \(d\in M_{xy}\).
The sets need not contain every forbidden parameter.

We also maintain a \emph{required-edge graph} \(J\) on \(X\): every edge
of \(J\) is known to be an edge of \(G[X]\).  Edges omitted from \(J\)
are unspecified.  Initially \(X=S\), and \(J\) consists of the clique
edges and the steps of the expanded paths in
Lemma~\ref{finite:labels}.  Thus neither the pattern nor \(J\) asserts
that an additional edge of \(G[S]\) is absent.

\begin{lemma}[Extension rule]\label{finite:extension-rule}
Let \(\mathcal P\) be the chosen optimal path system, with total amount
\(L<h=k+1-t\) and \(e\) paths, labelled as in
Lemma~\ref{finite:labels}.  Fix \(x<y\) in \(S\), and choose
\(z\in T(x)\), \(w\in T(y)\) with \(z<w\).  Let
\(E\subseteq\mathcal E(P)\) consist of old endpoint edges whose
represented paths do not contain \(x\) or \(y\) in their interiors.
Suppose that \(z,w\) have degree at most one in \(E\) and belong to
different components of the forest \(E\), with absent vertices treated
as isolated.  Here \(V(E)\) denotes the set of endpoints of edges in
\(E\), excluding all isolated clique vertices. Put
\[
 \begin{split}
 e'&=|E|+1,\\
 v'&=\bigl|V(E)\cup\{z,w\}\bigr|,\\
 q&=\sum_{(a,b)\in E}(b-a-1)+|x-z|+|y-w|.
 \end{split}
\]
Every positive integer \(d\) satisfying
\begin{equation}\label{finite:extension-interval}
 L+\mathbf 1_{\{e'\ge e\}}-q
 \ \le d\le\
 k+1-v'-q
\end{equation}
is a forbidden outside-path parameter between \(x,y\), relative to \(S\).
\end{lemma}

\begin{proof}
Suppose such an outside path exists.  Extend its end \(x\) to \(z\)
along the old path containing \(x\), or use the singleton segment
\(\{x\}\) if \(x\in Q\).  Extend its other end from \(y\) to \(w\)
in the same way.  Each extension meets \(Q\) only at its chosen
endpoint.

We first check that these two extension segments are disjoint.
If both are nontrivial and supported on the same old path with endpoints
\(a<b\), the conditions \(x<y\) and \(z<w\) force
\(z=a,w=b\).  The two segments then have vertex sets
\(\{a,\ldots,x\}\) and \(\{y,\ldots,b\}\), which are disjoint.
If their supporting old paths are different, their interiors are
disjoint, so an intersection could only be their common chosen clique
endpoint.  This would give \(z=w\), contrary to the choice.
If one segment is a singleton in \(Q\), it likewise can meet the other
segment only at that segment's chosen clique endpoint, again excluded.
Two singleton segments are disjoint because \(x<y\).

The assumed outside path has all its internal vertices outside \(S\).
Concatenating it with the two disjoint extension segments therefore
gives a simple \(z\)-\(w\) path whose interior lies outside \(Q\).
It has
\[
 d+|x-z|+|y-w|
\]
internal vertices: when \(x\notin Q\), the first extension contributes
all its vertices except \(z\), including \(x\), and similarly at \(y\).
In particular its amount is positive.

Every old path containing \(x\) or \(y\) in its interior must be
discarded, because that vertex is internal to the new path.  These
are exactly the old paths whose interiors can meet an extension.
If \(x,y\) lie on the same old path it is discarded only once.
Consequently the new path and all the old paths represented by \(E\)
have pairwise disjoint interiors.  Taking a subset of the remaining
old paths only discards more paths and preserves this property.

The endpoint graph \(E\) is already a linear forest.  Adding \(zw\)
keeps every degree at most two, by the degree assumptions, and creates
no cycle, because \(z,w\) were in different components.  This condition
also prevents a repeated endpoint edge.  Hence the old paths
represented by \(E\), together with the new path, form an admissible
system with \(e'\) paths, \(v'\) incident clique vertices, and amount
\[
 A=q+d.
\]
The lower bound in \eqref{finite:extension-interval} says that
\(A\ge L+1\) when \(e'\ge e\), and \(A\ge L\) when \(e'<e\).
If \(A<h\), it contradicts maximality of \(L\), or, in the case
\(A=L\) and \(e'<e\), the minimum-path tie-break.
If \(A\ge h\), the upper bound gives
\[
 h\le A\le k+1-v',
\]
contrary to Lemma~\ref{path:interval}.
Since \(L<h\), these alternatives cover the entire stated interval.
\end{proof}

For fixed \(x<y\), we put into \(M_{xy}\) the union of the positive
integer intervals \eqref{finite:extension-interval} over all allowed
choices of \(z,w,E\), and set \(M_{yx}=M_{xy}\).
Empty intervals contribute nothing.  This constructs the initial
constraints with empty diagonal.

For example, when \(k=5,t=3\) and \(P=((),(2))\), we have
\(Q=\{0,1,4\}\) and the old path \(1,2,3,4\).
Choose \(x=0,y=2,z=0,w=4\) and \(E=\varnothing\).
Then \(L=2\), \(e=e'=1\), \(v'=2\) and \(q=2\), so
\eqref{finite:extension-interval} forbids \(d=1,2\).
An outside \(0\)--\(2\) path relative to \(S\) with parameter \(d\)
would extend through \(2,3,4\); for \(d=1\) it closes through the
unused clique vertex \(1\), and for \(d=2\) it closes along the
clique edge \(4\)--\(0\).  Both give a cycle of order six.

The forest condition in Lemma~\ref{finite:extension-rule} also has a
simple interpretation under the labelling of
Lemma~\ref{finite:labels}.  Each component of \(E\) is a subchain of
an original monotonically labelled chain.  Thus, for \(z<w\),
connectivity from \(z\) to \(w\) is equivalent to the existence of a
route using only increasing old edges.  Such a route is found by
scanning the old edges in increasing chain order, starting with the
reachable set \(\{z\}\) and adjoining \(b\) whenever an edge
\((a,b)\in E\) has \(a\) already reachable.  The degree and connectivity
tests therefore implement precisely the forest condition used above.

\begin{lemma}[Required-path rule]\label{finite:required-path-rule}
Let \(X\supseteq S\), let \(J\) be a required-edge graph on \(X\), and
let \(x,y\in X\) be distinct.  If \(J\) contains a simple
\(x\)-\(y\) path of length \(\ell\), where \(1\le\ell\le k\), then
\[
 k-\ell\in M_{xy}
\]
is a valid forbidden outside-path constraint relative to \(X\).
\end{lemma}

\begin{proof}
If an outside path with \(d=k-\ell\) internal vertices existed, it
would have length \(d+1\).  Its interior is disjoint from the required
path, since the latter lies wholly in \(X\).  Their union is therefore
a simple cycle of length
\[
 \ell+(d+1)=k+1.
\]
For \(d=0\), the outside path is the edge \(xy\) and the required path
has length \(k\ge5\), so the same conclusion holds without any repeated
edge.  Additional edges of \(G\) may be chords of this cycle and do not
affect the contradiction.
\end{proof}

Applying Lemma~\ref{finite:required-path-rule} to all simple paths in
\(J\) preserves symmetry, because each path can be reversed.
It adds no diagonal entries, since a positive-length simple path has
distinct endpoints.  A required edge \(xy\in E(J)\) with
\(0\in M_{xy}\) is consequently an immediate contradiction.

\begin{lemma}[Inheritance under enlargement]\label{finite:inheritance}
Suppose \(S\subseteq X\subseteq X'\).  Every forbidden outside-path
constraint relative to \(X\), for endpoints in \(X\), remains valid
relative to \(X'\).
\end{lemma}

\begin{proof}
An outside path relative to \(X'\) has all its internal vertices in
\(V(G)\setminus X'\subseteq V(G)\setminus X\), so it is also an outside
path relative to \(X\).  This includes parameter \(0\), for which the
interior is empty.
\end{proof}

In particular, when a hypothesis introduces a fresh vertex into \(X\),
the old constraints may be retained for all old endpoint pairs.
Pairs involving the new vertex begin with no inherited constraints;
Lemma~\ref{finite:required-path-rule} then provides exclusions using
the enlarged required-edge graph.  Every exclusion obtained so far is
therefore justified by an actual path construction, independently of
any unspecified edges of \(G\).
